% Informe del sprint 3 del framework de desarrollo seguro.
% Compilar con XeLaTeX (dos pasadas para el índice):
%   latexmk -xelatex informe_sprint3.tex
\documentclass[11pt,a4paper]{report}

% ---------- Idioma y tipografía ----------
\usepackage{fontspec}
\usepackage[spanish,es-noshorthands,es-tabla]{babel}
\setmonofont{Latin Modern Mono}[Scale=0.95]

% ---------- Página ----------
\usepackage[a4paper,top=2.6cm,bottom=2.6cm,left=2.5cm,right=2.5cm,headheight=14pt]{geometry}
\usepackage{xcolor}
\usepackage{titlesec}
\usepackage{fancyhdr}
\usepackage{booktabs}
\usepackage{tabularx}
\usepackage{longtable}
\usepackage{array}
\usepackage{colortbl}
\usepackage{enumitem}
\usepackage{caption}
\usepackage{float}
\usepackage[most]{tcolorbox}
\usepackage{tikz}
\usetikzlibrary{arrows.meta,positioning,fit,backgrounds}
\usepackage[htt]{hyphenat}  % permite cortar los nombres largos en \texttt
\usepackage{xurl}
\usepackage[hidelinks]{hyperref}

% ---------- Colores (los mismos del plan técnico v2) ----------
\definecolor{azulfw}{HTML}{2A64A8}
\definecolor{azulclaro}{HTML}{E4EEF9}
\definecolor{verdefw}{HTML}{3E6B12}
\definecolor{verdeclaro}{HTML}{EAF1E1}
\definecolor{ocrefw}{HTML}{8A4E0B}
\definecolor{ocreclaro}{HTML}{FCF0DF}
\definecolor{rojofw}{HTML}{A33A3A}
\definecolor{rojoclaro}{HTML}{F8E3E3}
\definecolor{grisfw}{HTML}{6B6B6B}

% ---------- Títulos ----------
\titleformat{\chapter}[hang]{\huge\bfseries\color{azulfw}}{\thechapter.}{0.5em}{}
\titlespacing*{\chapter}{0pt}{0pt}{1.2em}
\titleformat{\section}{\Large\bfseries\color{azulfw}}{\thesection}{0.6em}{}
\titleformat{\subsection}{\large\bfseries}{\thesubsection}{0.6em}{}
\addto\captionsspanish{\renewcommand{\contentsname}{\color{azulfw}Índice}}

% ---------- Encabezado y pie ----------
\pagestyle{fancy}
\fancyhf{}
\fancyhead[L]{\small\color{grisfw}Framework de desarrollo seguro guiado por conocimiento}
\fancyhead[R]{\small\color{grisfw}Informe del sprint 3}
\fancyfoot[C]{\small\thepage}
\renewcommand{\headrulewidth}{0.4pt}
\fancypagestyle{plain}{\fancyhf{}\fancyfoot[C]{\small\thepage}\renewcommand{\headrulewidth}{0pt}}

% ---------- Tablas ----------
\renewcommand{\arraystretch}{1.25}
\newcolumntype{L}{>{\raggedright\arraybackslash}X}
\newcolumntype{P}[1]{>{\raggedright\arraybackslash}p{#1}}
\newcolumntype{C}[1]{>{\centering\arraybackslash}p{#1}}
\captionsetup{font=small,labelsep=colon}
\setlist{itemsep=0.25em,topsep=0.4em}

% ---------- Recuadros ----------
\newtcolorbox{cajaverde}[1]{enhanced,colback=verdeclaro,colframe=verdefw,
  coltitle=white,fonttitle=\bfseries,title={#1},arc=1.5mm,boxrule=0.6pt,
  left=2mm,right=2mm,top=1mm,bottom=1mm}
\newtcolorbox{cajaazul}[1]{enhanced,colback=azulclaro,colframe=azulfw,
  coltitle=white,fonttitle=\bfseries,title={#1},arc=1.5mm,boxrule=0.6pt,
  left=2mm,right=2mm,top=1mm,bottom=1mm}
\newtcolorbox{cajaocre}[1]{enhanced,colback=ocreclaro,colframe=ocrefw,
  coltitle=white,fonttitle=\bfseries,title={#1},arc=1.5mm,boxrule=0.6pt,
  left=2mm,right=2mm,top=1mm,bottom=1mm}

\newcommand{\pr}[1]{\href{https://github.com/PracticasESPE/Framework_Desarrollo_Seguro/pull/#1}{\##1}}
\newcommand{\run}[1]{\href{https://github.com/PracticasESPE/Framework_Desarrollo_Seguro/actions/runs/#1}{#1}}
\newcommand{\filasprint}[1]{\rowcolor{azulclaro}\multicolumn{5}{l}{\bfseries #1}\\}

\hypersetup{pdftitle={Informe del sprint 3},pdfauthor={Marcelo Acuña Gamboa, Abner Arboleda}}

\begin{document}

% ======================================================================
% PORTADA
% ======================================================================
\begin{titlepage}
\centering
\vspace*{4.2cm}
{\color{azulfw}\rule{\textwidth}{0.8pt}}\par\vspace{0.9cm}
{\fontsize{24}{30}\selectfont\bfseries Framework de desarrollo seguro\\
basado en métodos computacionales\\ guiados por conocimiento\par}
\vspace{0.9cm}{\color{azulfw}\rule{\textwidth}{0.8pt}}\par
\vspace{1.1cm}
{\Large Informe del sprint 3\par}
\vspace{0.3cm}
{\large Fases 1 y 2: del documento de requisitos a la arquitectura validada\par}
\vspace{0.3cm}
{\normalsize 4 al 13 de noviembre de 2026\par}
\vspace{2.2cm}
{\normalsize
\begin{tabular}{rl}
Marcelo Acuña Gamboa & Participante 1: motor e IA\\
Abner Arboleda & Participante 2: plataforma y seguridad\\
\end{tabular}\par}
\vfill
{\normalsize Universidad de las Fuerzas Armadas – ESPE \quad$\bullet$\quad Ingeniería de Software\par}
\vspace{0.4cm}
{\normalsize Versión 1.0 \quad$\bullet$\quad 15 de noviembre de 2026\par}
\end{titlepage}

\tableofcontents

% ======================================================================
\chapter{Resumen del sprint}
% ======================================================================

\section{Objetivo}

El sprint 3 construye las dos primeras fases del framework. La fase 1 convierte el documento de
requisitos del cliente en un modelo estructurado y clasificado; la fase 2 propone la arquitectura
de microservicios, la valida con OPA y la documenta con ADR y un diagrama C4. Además, el sprint
debía entregar el dataset de evaluación y la primera versión del dashboard.

\begin{cajaazul}{Sin hito propio: prepara el del 27 de noviembre}
El sprint no tiene hito de control. Entrega la entrada del hito del 27/11 (\emph{el framework
genera microservicios desde un documento de requisitos y les aplica el Security Gate 1}):
\texttt{modelo\_requisitos.json} y \texttt{arquitectura.json}.
\end{cajaazul}

\begin{cajaverde}{Resultado: las fases 1 y 2 funcionan de punta a punta}
El 15/11/2026 se ejecutó la cadena completa con los tres documentos del dataset: carga,
extracción con Docling, clasificación con el LLM, registro en Neo4j, propuesta de arquitectura con
el ciclo LLM ↔ OPA, ADR y diagrama C4. Los tres documentos llegaron al final sin intervención
manual: se encontraron los 111 requisitos de la referencia, OPA aprobó las tres arquitecturas y
quedaron 29 decisiones enlazadas a sus requisitos y reglas en el grafo (capítulo~\ref{cap:prueba}).
\end{cajaverde}

\section{Indicadores}

\begin{table}[H]
\centering
\caption{Indicadores del sprint 3, medidos el 15/11/2026 sobre \texttt{dev} y en local.}
\begin{tabularx}{\textwidth}{@{}P{6.2cm}L@{}}
\toprule
\textbf{Indicador} & \textbf{Valor}\\
\midrule
Días y horas del sprint & 8 días, 64 horas de equipo (32 por integrante)\\
Pull requests fusionados en \texttt{dev} & 24 (\pr{179} a \pr{202}; 10 de Marcelo y 14 de Abner), más el de este informe\\
Duración del pipeline en CI & 12 min 20 s (ejecución \run{43660815603}, todas las fases en verde)\\
Requisitos encontrados (fase 1) & 111 de 111 en los tres documentos; 1 falso positivo\\
Clasificación funcional / no funcional & 86\,\% de aciertos (96 de 111)\\
Arquitecturas aprobadas por OPA & 3 de 3, las tres al primer intento; 29 nodos \texttt{Decision} en Neo4j\\
Agrupación de requisitos en servicios & F1 por pares de 0,82, 0,66 y 0,83 frente a la referencia\\
Pruebas del motor (pytest) & 676 en verde y 15 omitidas (requieren Neo4j o son las lentas de PDF); 94\,\% de cobertura de \texttt{orchestrator/}\\
Pruebas de políticas Rego & 87 en verde (gates y arquitectura)\\
Pruebas del dashboard (Vitest) & 33 en verde; ESLint con \texttt{eslint-plugin-security} y Semgrep en su propio job de CI\\
Alertas abiertas en GitHub Security & 47, las mismas del sprint 2, todas de severidad advertencia o baja\\
\bottomrule
\end{tabularx}
\end{table}

% ======================================================================
\chapter{Trabajo realizado}
% ======================================================================

\section{Cumplimiento del cronograma}

La tabla compara el cronograma del plan con lo entregado. La columna PR indica el pull request con
el que cada tarea entró en \texttt{dev}.

{\small
\begin{longtable}{@{}P{1.55cm}P{4.6cm}P{4.6cm}P{1.6cm}P{1.6cm}@{}}
\caption{Tareas del sprint 3 por integrante y pull request asociado.}\\
\toprule
\textbf{Fecha} & \textbf{Marcelo (motor e IA)} & \textbf{Abner (plataforma y seguridad)} & \textbf{PR Marcelo} & \textbf{PR Abner}\\
\midrule
\endfirsthead
\toprule
\textbf{Fecha} & \textbf{Marcelo (motor e IA)} & \textbf{Abner (plataforma y seguridad)} & \textbf{PR Marcelo} & \textbf{PR Abner}\\
\midrule
\endhead
\filasprint{Sprint 3: fases 1 y 2}
04/11 mié & Docling: extracción de PDF, DOCX y TXT. & Dataset: documento de requisitos 1 (tienda en línea). & \pr{185} & \pr{179}\\
05/11 jue & Normalización del texto y esquema Pydantic \texttt{Requisito}. & Documentos 2 y 3 (gestión académica y citas médicas). & \pr{186} & \pr{180}\\
06/11 vie & Clasificación RF/RNF y etiquetas de seguridad con LLM. & Dashboard React + Vite: estructura y carga del documento. & \pr{187}, \pr{188} & \pr{182}, \pr{184}\\
09/11 lun & Requisitos guardados como nodos en Neo4j. & Dashboard: progreso por fase con Server-Sent Events. & \pr{192} & \pr{190}\\
10/11 mar & Propuesta de arquitectura con LLM. & API del motor: endpoints de ejecución y estado. & \pr{194} & \pr{191}\\
11/11 mié & OPA: reglas arquitectónicas e inferencia de patrones. & Generación de ADR en Markdown. & \pr{196} & \pr{197}\\
12/11 jue & Ciclo LLM ↔ OPA y nodos \texttt{Decision}. & Diagrama C4 en Mermaid. & \pr{202} & \pr{201}\\
13/11 vie & Fases 1--2 con los 3 documentos y demostración. & Revisión de ADR y diagramas, correcciones del dashboard e informe. & pendiente & este PR\\
\bottomrule
\end{longtable}}

Además del cronograma entraron trabajos que surgieron durante el sprint: dos excepciones con
caducidad en el Gate 3 por CVE de \texttt{zlib} en la imagen base (\pr{181}, \pr{183}), la
corrección de la lectura de PDF en rutas con tildes (\pr{189}), la gestión de la cuota del LLM con
caché y respaldo entre proveedores (\pr{193}), el registro automático de los requisitos en Neo4j
desde el dashboard (\pr{198}), la segmentación de documentos sin estructura (\pr{200}) y un script
de demostración del ciclo LLM ↔ OPA que sirvió de guía para \pr{202} (\pr{199}).

Las tareas de plataforma entraron en su fecha o con un día de diferencia (\pr{180} y \pr{182}).
En el motor, \pr{192}, \pr{194} y \pr{196} entraron a tiempo, pero las de la fase 1 se retrasaron
entre tres y cuatro días (\pr{185} el 07/11, \pr{186} y \pr{187} el 09/11) y el ciclo LLM ↔ OPA
tres días (\pr{202} el 15/11). Al cierre, la issue \#58 sigue abierta aunque su PR está fusionado,
porque el PR no incluía \texttt{closes}, y la tarea \#60 está pendiente.

\section{Motor e IA}

\subsection{Fase 1: del documento al modelo de requisitos}
\begin{itemize}
\item \textbf{Extracción} (\pr{185}): Docling convierte PDF y DOCX a Markdown sin OCR, con límite
de páginas y de tiempo; TXT y Markdown se leen directamente. Las pruebas comprueban que se
conservan todas las frases de la referencia en los tres documentos y los cuatro formatos.
\item \textbf{Normalización y esquema} (\pr{186}, \pr{188}, \pr{200}): el texto se limpia de
caracteres invisibles y se divide en unidades con su ubicación (sección, viñeta o párrafo). Los
documentos sin títulos se dividen en oraciones.
\item \textbf{Clasificación} (\pr{187}): el LLM clasifica cada sección en un lote. El texto del
cliente va entre marcas con un código aleatorio para que no pueda cerrar el bloque ni dar órdenes
al modelo, y la respuesta se valida con Pydantic. El resultado es
\texttt{modelo\_requisitos.json}, con su JSON Schema en \texttt{docs/esquemas/}.
\item \textbf{Registro en Neo4j} (\pr{192}): \texttt{(:Ejecucion)-[:TIENE\_REQUISITO]->(:Requisito)}.
\end{itemize}

\subsection{Fase 2: de los requisitos a la arquitectura}
\begin{itemize}
\item \textbf{Propuesta} (\pr{194}): el modelo de razonamiento propone bounded contexts,
microservicios con su base de datos, dependencias síncronas o por eventos y decisiones con su
justificación. El motor comprueba que cada requisito quede asignado una sola vez.
\item \textbf{Reglas de OPA} (\pr{196}): cuatro reglas bloqueantes (base de datos propia,
dependencias conocidas, servicios únicos y sin ciclos de llamadas síncronas), dos advertencias y la
inferencia de los patrones de la base de conocimiento a partir de los rasgos de la arquitectura.
\item \textbf{Ciclo LLM ↔ OPA} (\pr{202}): si OPA rechaza, el LLM recibe los motivos de todos los
rechazos anteriores, como mucho tres veces. Cada intento queda en
\texttt{<nombre>.intentos.json} y las decisiones aprobadas en Neo4j como nodos \texttt{Decision}
enlazados a sus requisitos (\texttt{ORIGINA}), a las reglas (\texttt{BASADA\_EN}) y a los patrones
(\texttt{APLICA}). En este PR la regla de ciclos pasó a mirar solo las llamadas síncronas: un ciclo
de eventos es una saga coreografiada y ahora genera una advertencia.
\end{itemize}

\section{Plataforma y seguridad}

\subsection{Dataset de evaluación}
Tres documentos de requisitos (\pr{179}, \pr{180}) escritos como los entregaría un cliente:
tienda en línea, gestión académica y citas médicas, cada uno en PDF, DOCX, TXT y Markdown. Cada
documento tiene su referencia (\texttt{referencia.yaml}) con 30, 38 y 43 requisitos clasificados,
sus etiquetas de seguridad y los servicios esperados. Es también la entrada de los experimentos
DSR, por lo que queda congelado.

\subsection{Dashboard y API del motor}
\begin{itemize}
\item \textbf{Dashboard v0 y v1} (\pr{182}, \pr{184}, \pr{190}): React + Vite, servido por la
propia API sin CORS. Carga del documento o escritura de los requisitos, progreso de cada fase en
tiempo real con Server-Sent Events, cancelación, reanudación al recargar y tabla de los requisitos
identificados.
\item \textbf{API} (\pr{191}): \texttt{POST /documentos} con validación del contenido;
\texttt{POST /ejecuciones}, estado, eventos SSE y cancelación; \texttt{GET
/documentos/\{id\}/requisitos}. Documentada en \texttt{/docs}.
\item \textbf{Registro automático en Neo4j} (\pr{198}): al terminar la fase 1 desde el dashboard,
los requisitos quedan en el grafo sin ejecutar comandos; si Neo4j no está, la fase termina igual y
lo avisa.
\end{itemize}

\subsection{Gestión de los proveedores de LLM}
La cuota gratuita de Groq se agotó durante las pruebas del 09/11. \pr{193} añadió:
\begin{itemize}
\item una caché en disco de las respuestas (repetir un documento no gasta cuota);
\item el tratamiento del límite de uso (se respeta \texttt{retry-after} y se distingue la cuota
diaria);
\item el respaldo entre proveedores (\texttt{LLM\_RESPALDO}) y la medición de tokens por documento.
\end{itemize}

\begin{table}[H]
\centering
\caption{Groq frente a Gemini con los tres documentos (10/11/2026, caché vacía).}
\begin{tabularx}{\textwidth}{@{}LC{2.0cm}C{1.6cm}C{2.2cm}C{2.4cm}@{}}
\toprule
\textbf{Proveedor} & \textbf{Tokens} & \textbf{Tiempo} & \textbf{Tipo correcto} & \textbf{F1 de etiquetas}\\
\midrule
Groq (\texttt{gpt-oss-20b}) & 96\,167 & 589 s & 89\,\% & 0,67\\
Gemini (\texttt{gemini-3.5-flash-lite}) & 45\,011 & 142 s & 85\,\% & 0,66\\
\bottomrule
\end{tabularx}
\end{table}

Groq sigue como proveedor principal y Gemini como respaldo. Para la evaluación de la tesis se usa un
solo proveedor y sin respaldo, porque mezclar modelos cambia los resultados.

\subsection{ADR y diagrama C4}
\begin{itemize}
\item \textbf{ADR} (\pr{197}): un ADR por decisión de la arquitectura y otro por patrón que infiere
OPA, solo si OPA aprobó. Los hechos (decisión, requisitos, controles ASVS y CWE con su nombre) salen
de datos validados; el LLM rápido solo redacta el contexto y las consecuencias, como indica la
tabla 6.1 del plan, y su texto se descarta si cita identificadores que no son del ADR.
\item \textbf{Diagrama C4} (\pr{201}): \texttt{docs/arquitectura.md} con el contexto (actores de
los requisitos) y los contenedores (gateway, identidad, bus de eventos, servicios y sus bases de
datos), sin LLM. El nivel de contenedores se dibuja con un \texttt{flowchart} con el estilo de C4
porque la sintaxis \texttt{C4Container} de Mermaid superponía las flechas.
\end{itemize}

% ======================================================================
\chapter{Prueba de las fases 1 y 2 con los tres documentos}\label{cap:prueba}
% ======================================================================

El 15/11 se ejecutó la cadena completa con la versión DOCX de cada documento, con Groq como
proveedor y Gemini como respaldo, Neo4j levantado y la caché vacía para la fase 1:

\begin{center}
\small\ttfamily
requisitos clasificar → evaluar (fase 1) → arquitectura proponer → evaluar (fase 2)\\
→ arquitectura adr → arquitectura diagrama
\end{center}

\section{Fase 1: requisitos}

\begin{table}[H]
\centering
\caption{Fase 1 frente a la referencia del dataset.}
\begin{tabularx}{\textwidth}{@{}LC{1.9cm}C{1.6cm}C{1.9cm}C{3.4cm}@{}}
\toprule
\textbf{Documento} & \textbf{Encontrados} & \textbf{Falsos positivos} & \textbf{Tipo correcto} & \textbf{Etiquetas (P / R / F1)}\\
\midrule
1 Tienda en línea & 30/30 & 0 & 87\,\% & 100\,\% / 53\,\% / 0,69\\
2 Gestión académica & 38/38 & 0 & 87\,\% & 100\,\% / 39\,\% / 0,56\\
3 Citas médicas & 43/43 & 1 & 86\,\% & 70\,\% / 36\,\% / 0,47\\
\midrule
\textbf{Total} & \textbf{111/111} & \textbf{1} & \textbf{86\,\%} & F1 medio 0,57\\
\bottomrule
\end{tabularx}
\end{table}

La extracción y la segmentación no pierden requisitos, y la distinción entre funcional y no
funcional acierta en 86\,\% de los casos. El punto débil son las \textbf{etiquetas de seguridad}:
la precisión es alta, pero el modelo omite entre la mitad y dos tercios de las etiquetas que
esperaba la referencia. Como las etiquetas deciden qué patrones y controles infiere OPA, mejorar su
exhaustividad es la prioridad de la fase 1. Los resultados varían entre ejecuciones: el mismo
documento 1 dio un F1 de 0,76 el 10/11.

\section{Fase 2: arquitectura}

\begin{table}[H]
\centering
\caption{Fase 2: ciclo LLM ↔ OPA y comparación con los servicios de la referencia.}
\begin{tabularx}{\textwidth}{@{}LC{1.7cm}C{1.9cm}C{2.4cm}C{2.4cm}@{}}
\toprule
\textbf{Documento} & \textbf{Intentos} & \textbf{Servicios (ref.)} & \textbf{Nombres que coinciden} & \textbf{F1 de agrupación}\\
\midrule
1 Tienda en línea & 1 & 6 (5) & 3 & 0,82\\
2 Gestión académica & 1 & 6 (6) & 2 & 0,66\\
3 Citas médicas & 1 & 7 (6) & 5 & 0,83\\
\bottomrule
\end{tabularx}
\end{table}

OPA aprobó las tres propuestas al primer intento e infirió en las tres los cinco patrones de la
base de conocimiento que aplican (API Gateway, Database per Service, Event-Driven, Circuit Breaker
y Saga). En Neo4j quedaron 10, 10 y 9 nodos \texttt{Decision}. La agrupación de requisitos en
servicios es buena en los documentos 1 y 3. En el documento 2 el LLM organizó los servicios de
otra forma (por ejemplo, \texttt{auth} e \texttt{integraciones} en lugar de \texttt{usuarios} y
\texttt{notificaciones}), lo que baja el F1 sin que la arquitectura sea incorrecta. El F1 de
agrupación mide la coincidencia con una propuesta de referencia, no la validez de la arquitectura.

Antes de \pr{202}, la regla de ciclos también bloqueaba los ciclos que pasan por eventos. En
pruebas del 10/11 y 11/11 OPA rechazó por ese motivo propuestas del LLM, y al devolverle los
motivos se aprobaron en el segundo o tercer intento.

\section{ADR, diagrama y coste}

\begin{table}[H]
\centering
\caption{Salidas finales y coste de cada documento (15/11/2026).}
\begin{tabularx}{\textwidth}{@{}LC{2.4cm}C{2.4cm}C{2.4cm}C{1.7cm}@{}}
\toprule
\textbf{Documento} & \textbf{ADR (plantilla)} & \textbf{Tokens fase 1} & \textbf{Tokens ADR} & \textbf{Tiempo}\\
\midrule
1 Tienda en línea & 10 (0) & 26\,220 & 23\,996 & 341 s\\
2 Gestión académica & 10 (1) & 31\,917 & 25\,128 & 402 s\\
3 Citas médicas & 9 (0) & 33\,464 & 20\,757 & 441 s\\
\bottomrule
\end{tabularx}
\end{table}

Un ADR del documento 2 se redactó con la plantilla del motor porque el LLM no cumplió el esquema en
tres intentos; el ADR lo indica en su campo «Redacción». Los tres diagramas C4 se generaron con
todos sus patrones enlazados a su ADR. Con la cuota gratuita de Groq (200\,000 tokens al día), el
ciclo completo de un documento cuesta alrededor de un tercio de la cuota diaria; con la caché,
repetirlo no cuesta nada.

% ======================================================================
\chapter{Revisión de los ADR y diagramas, y correcciones del dashboard}
% ======================================================================

La revisión de la tarea \#59 regeneró los ADR y el diagrama del documento 1 con las reglas de
\pr{202}, comparó el resultado con lo que estaba en el repositorio y recorrió el dashboard con el
flujo completo en escritorio y en un móvil de 390 píxeles.

\section{Hallazgos y correcciones}

{\small
\begin{longtable}{@{}P{3.6cm}P{5.2cm}P{5.4cm}@{}}
\caption{Hallazgos de la revisión del sprint 3 y su corrección (este PR).}\\
\toprule
\textbf{Hallazgo} & \textbf{Causa} & \textbf{Corrección}\\
\midrule
\endfirsthead
\toprule
\textbf{Hallazgo} & \textbf{Causa} & \textbf{Corrección}\\
\midrule
\endhead
El ADR-0001 del repositorio citaba las marcas internas del prompt (\texttt{<requisito-f845ef5a …>}) & El LLM copió el formato de los datos en lugar de redactar. Ocurre de forma intermitente: en la regeneración volvió a pasar en otro ADR. No es un riesgo de seguridad, porque el texto está escapado. & El prompt pide citar los requisitos solo por su id; el motor rechaza una redacción con esas marcas y la pide otra vez indicando el motivo; si vuelve a fallar, usa la plantilla. Dos pruebas nuevas.\\
El índice de los ADR y \texttt{docs/arquitectura.md} describían la regla de ciclos anterior & \pr{202} cambió la regla a ciclos de llamadas síncronas después de generarlos. & Se regeneraron \texttt{docs/adr/} y \texttt{docs/arquitectura.md}.\\
El diagrama de contexto mostraba actores duplicados o que no son personas & El LLM escribe el actor de cada requisito libremente: \emph{medico} y \emph{médico}, \emph{perfil} (la tabla de roles) o menciones sueltas como \emph{empresa}. & Se unen las variantes con y sin tilde, se excluye \emph{perfil} y los actores de un solo requisito se agrupan en «Otros actores». Cuatro pruebas nuevas.\\
En el móvil la página se desbordaba 201 píxeles hacia la derecha & Tres causas: la columna de la tarjeta de requisitos crecía hasta el ancho de su tabla; las rutas largas en \texttt{<code>} no se partían; y el texto oculto «Acción» de la tabla de documentos, con posición absoluta, escapaba de su contenedor. & \texttt{minmax(0, 1fr)} en las tarjetas, \texttt{overflow-wrap} en el código y \texttt{position: relative} en el contenedor de las tablas. Desborde medido después: 0 píxeles.\\
El aviso de API apagada indicaba \texttt{uv run framework api} & Ese comando no arranca en Windows cuando la ruta del proyecto tiene tildes. & El aviso indica \texttt{uv run python -m orchestrator.cli api}, que funciona en todos los casos.\\
\bottomrule
\end{longtable}}

\section{Lo que queda para el motor}

\begin{itemize}
\item Los actores salen con nombres poco cuidados (\emph{direccionmedica}, \emph{coordinacion} sin
tilde). Conviene normalizarlos en la clasificación de la fase 1, no en el diagrama.
\item La exhaustividad de las etiquetas de seguridad (36 a 53\,\%) limita lo que OPA puede inferir.
\item El ADR que se redactó con plantilla en el documento 2 muestra que un esquema más simple para la
redacción reduciría los reintentos.
\end{itemize}

% ======================================================================
\chapter{Estado del pipeline al cierre}
% ======================================================================

\section{Las fases 1 y 2}

\begin{figure}[H]
\centering
\begin{tikzpicture}[
  font=\sffamily\footnotesize,
  caja/.style={draw=azulfw,fill=azulclaro,rounded corners=2pt,align=center,minimum height=1.05cm,text width=2.25cm},
  gate/.style={draw=ocrefw,fill=ocreclaro,rounded corners=2pt,align=center,minimum height=1.05cm,text width=2.25cm},
  fin/.style={draw=verdefw,fill=verdeclaro,rounded corners=2pt,align=center,minimum height=0.95cm,text width=3.0cm},
  flecha/.style={-{Stealth[length=2.2mm]},thick,color=grisfw},
  node distance=0.55cm and 0.45cm]
\node[caja] (doc) {Documento\\(dashboard o API)};
\node[caja,right=of doc] (f1) {Fase 1\\Docling y\\clasificación LLM};
\node[caja,right=of f1] (llm) {Fase 2\\propuesta\\del LLM};
\node[gate,right=of llm] (opa) {OPA\\4 reglas y\\patrones};
\node[fin,below=1.0cm of f1] (req) {Neo4j: Requisito};
\node[fin,below=1.0cm of opa] (dec) {Neo4j: Decision\\→ Regla, Patrón};
\node[fin,right=1.0cm of dec] (adr) {ADR y C4\\(\texttt{docs/})};
\draw[flecha] (doc) -- (f1);
\draw[flecha] (f1) -- (llm);
\draw[flecha] (llm) -- (opa);
\draw[flecha] (opa.north) to[bend right=35] node[above,font=\sffamily\scriptsize]{rechazo: motivos (máx. 3)} (llm.north);
\draw[flecha] (f1) -- (req);
\draw[flecha] (opa) -- (dec);
\draw[flecha] (dec) -- (adr);
\end{tikzpicture}
\caption{Recorrido de las fases 1 y 2. Si OPA rechaza la propuesta, el LLM recibe los motivos y
propone otra; solo una arquitectura aprobada genera ADR y diagrama.}
\end{figure}

\section{Fases del pipeline}

\begin{table}[H]
\centering
\caption{Estado de cada fase al cierre del sprint 3.}
\begin{tabularx}{\textwidth}{@{}P{0.9cm}P{4.0cm}P{2.6cm}L@{}}
\toprule
\textbf{Fase} & \textbf{Nombre} & \textbf{Estado} & \textbf{Detalle}\\
\midrule
1 & Adquisición de requisitos & Completa en local & Dashboard, API y CLI. En CI solo se ejecuta la extracción con Docling del PDF del documento 1, sin LLM.\\
2 & Análisis y propuesta arquitectónica & Completa en local & CLI con ciclo LLM ↔ OPA, ADR y C4. En CI aparece como omitida.\\
3 & Generación & Versión inicial & Compila el servicio de referencia; la generación desde la arquitectura es el sprint 4.\\
4 & Security Gate 1 & Con decisión OPA & Sin cambios en este sprint.\\
5--6 & Refinamiento iterativo & Pendiente & Sprint 5.\\
7 & Integración & Parcial & Keycloak con guard JWT; Traefik y RabbitMQ en el sprint 5.\\
8--11 & Gates 2 y 3, pruebas e imagen & Con decisión OPA & Sin cambios salvo las excepciones de \texttt{zlib} del Gate 3.\\
\bottomrule
\end{tabularx}
\end{table}

Las fases 1 y 2 no se ejecutan completas en CI porque necesitan la clave del LLM y consumen cuota.
La planificación del sprint 3 proponía usar en los PR un \texttt{modelo\_requisitos.json} y un
\texttt{arquitectura.json} fijos y llamar al LLM solo en los \emph{push} a \texttt{dev}; queda
pendiente.

% ======================================================================
\chapter{Incidencias y lecciones aprendidas}\label{cap:incidencias}
% ======================================================================

Como en los sprints anteriores, ninguna incidencia se resolvió relajando un gate o una regla.

{\small
\begin{longtable}{@{}P{3.3cm}P{4.5cm}P{5.0cm}P{1.2cm}@{}}
\caption{Incidencias del sprint 3 y su corrección.}\\
\toprule
\textbf{Incidencia} & \textbf{Causa} & \textbf{Corrección} & \textbf{PR}\\
\midrule
\endfirsthead
\toprule
\textbf{Incidencia} & \textbf{Causa} & \textbf{Corrección} & \textbf{PR}\\
\midrule
\endhead
Gate 3 rechazado el 06/11 y el 07/11 & Dos CVE altas de \texttt{zlib} publicadas en la imagen \emph{distroless}, sin versión corregida todavía. & Excepciones en \texttt{.trivyignore.yaml} con justificación (Node no usa la \texttt{zlib} del sistema) y caducidad el 21/11. & \pr{181}, \pr{183}\\
La fase 1 fallaba con todos los PDF en Windows & docling-parse no abre sus propios recursos si la ruta tiene tildes. & El PDF se lee con pypdfium2. El CI no lo detectaba porque sus rutas son ASCII. & \pr{189}\\
Cuota diaria de Groq agotada el 09/11 & Unos 28\,000 tokens por documento frente a 200\,000 diarios del plan gratuito. & Caché de respuestas, \texttt{retry-after}, distinción de la cuota diaria y respaldo con Gemini. & \pr{193}\\
Gemini no podía clasificar & El adaptador enviaba el esquema en un formato que la API rechazaba, y Gemini devolvía las ubicaciones con entidades HTML. & JSON Schema estándar sin \texttt{maxItems}; las ubicaciones se decodifican. & \pr{193}\\
Una clave de API apareció en un mensaje de error & Se pegó en \texttt{LLM\_RESPALDO} y el error repetía el valor recibido. & Ningún error ni \texttt{/configuracion} repiten el valor; la CLI no muestra variables locales. Se recomendó rotar la clave. & \pr{193}\\
\texttt{dev} con marcas de conflicto & Se fusionó un PR con un conflicto sin resolver en \texttt{cli.py}; el CI falló. & PR de corrección inmediato. & \pr{195}\\
Repetir \texttt{arquitectura proponer} daba la misma propuesta rechazada & La caché devuelve la misma respuesta a la misma pregunta. & El ciclo pasa al LLM los motivos del rechazo, lo que cambia la pregunta. & \pr{199}, \pr{202}\\
Errores de TLS al descargar modelos (local) & El antivirus intercepta las conexiones TLS e inyecta su script en las páginas, que la CSP bloquea. & Diagnóstico; no requiere cambios en el código. & --\\
\bottomrule
\end{longtable}}

\begin{cajaocre}{Lecciones principales}
\begin{itemize}[leftmargin=1.2em]
\item \textbf{El LLM es un recurso limitado.} La cuota diaria condiciona cuánto se puede probar. La
caché y el respaldo son parte del diseño, no un añadido, y la evaluación debe fijar un único
proveedor.
\item \textbf{Validar lo que escribe el LLM, también la forma.} Además del esquema, hubo que
comprobar que no inventara identificadores ni copiara las marcas del prompt. Cada comprobación que
se añadió vino de un fallo observado en una ejecución real.
\item \textbf{Probar en el entorno del usuario.} Dos fallos (PDF con tildes y el comando
\texttt{framework}) solo aparecían en Windows con rutas no ASCII, que el CI no reproduce.
\item \textbf{Cerrar las issues con el PR.} Cuatro tareas fusionadas (\#46, \#48, \#50 y \#58)
quedaron con su issue abierta porque el PR no incluía \texttt{closes}; hubo que cerrarlas a mano.
\end{itemize}
\end{cajaocre}

% ======================================================================
\chapter{Seguridad del propio framework}
% ======================================================================

Controles del capítulo 10 del plan añadidos en el sprint 3:

{\small
\begin{longtable}{@{}P{4.2cm}P{10.2cm}@{}}
\caption{Controles implementados en el sprint 3.}\\
\toprule
\textbf{Riesgo} & \textbf{Control activo}\\
\midrule
\endfirsthead
\toprule
\textbf{Riesgo} & \textbf{Control activo}\\
\midrule
\endhead
Archivo malicioso como documento (CWE-434, 22, 400, 409) & Lista permitida de formatos y comprobación del contenido, máximo de 10 MB cortando la lectura, nombre generado por el servidor y límites de descompresión del DOCX.\\
Inyección de instrucciones en el LLM & El texto del cliente va entre marcas con un código aleatorio; toda respuesta se valida con Pydantic; el motor rechaza identificadores inventados y marcas copiadas.\\
Decisiones de arquitectura delegadas al LLM & OPA aprueba o rechaza con reglas versionadas y probadas; solo una arquitectura aprobada genera ADR y diagrama.\\
Contenido no confiable en Markdown y Mermaid & El texto de los requisitos y del LLM se escapa antes de escribirse: no puede abrir HTML, imágenes, enlaces ni instrucciones de Mermaid.\\
API sin autenticación & Escucha solo en 127.0.0.1; cabeceras CSP, \texttt{nosniff} y \texttt{frame-ancestors 'none'}; una ejecución a la vez y errores sin detalles internos.\\
Exposición de claves de LLM & Las claves se leen como \texttt{SecretStr}; ningún error ni endpoint repite valores de configuración; la caché no guarda prompts ni documentos.\\
Dependencias del dashboard & \texttt{npm ci -{}-ignore-scripts}, ESLint con \texttt{eslint-plugin-security} y Semgrep en un job propio de CI.\\
Vulnerabilidades sin corrección en la imagen base & Excepciones con justificación y fecha de caducidad; al caducar, el Gate 3 vuelve a bloquear.\\
\bottomrule
\end{longtable}}

% ======================================================================
\chapter{Pendientes y planificación del sprint 4}
% ======================================================================

\section{Pendientes}

\begin{table}[H]
\centering
\caption{Pendientes al cierre del sprint 3.}
\begin{tabularx}{\textwidth}{@{}LP{2.6cm}P{2.2cm}@{}}
\toprule
\textbf{Pendiente} & \textbf{Responsable} & \textbf{Para el}\\
\midrule
Tarea \#60: capturas de las fases 1 y 2 con los tres documentos (este informe adelanta las cifras) & Marcelo & 17/11\\
Cerrar la issue \#58, ya fusionada en \pr{202} & Marcelo & 17/11\\
El ruleset de \texttt{dev} está desactivado: no exige PR ni el check \texttt{Pipeline por fases} (pendiente del sprint 2) & Abner & 17/11\\
Activar las alertas de Dependabot (pendiente del sprint 2) & Abner & 17/11\\
Ejecutar las fases 1 y 2 en CI con entradas fijas en los PR y LLM solo en \texttt{dev} & Ambos & 20/11\\
Promover \texttt{dev} a \texttt{test} y \texttt{main}: llevan 64 commits de retraso & Ambos & antes del 27/11\\
Retirar las excepciones de \texttt{zlib} cuando se publique la imagen corregida (caducan el 21/11) & Abner & 21/11\\
Mejorar la exhaustividad de las etiquetas de seguridad y normalizar los actores & Marcelo & sprint 4\\
Unificar quién decide cada gate (script u OPA) y conservar los resultados de Neo4j de CI & Ambos & antes de la evaluación\\
\bottomrule
\end{tabularx}
\end{table}

\section{Objetivo del sprint 4}

\begin{cajaazul}{Sprint 4: fases 3 y 4 (16 al 27 de noviembre de 2026)}
Generar los microservicios NestJS a partir de \texttt{arquitectura.json} con una plantilla Copier
y el LLM, verificar que compilan y aplicarles el Security Gate 1. Hito del 27/11: \emph{fases 1-4,
genera microservicios y aplica Gate 1}.
\end{cajaazul}

{\small
\begin{longtable}{@{}P{1.55cm}P{5.6cm}P{5.6cm}@{}}
\caption{Cronograma del sprint 4 por integrante (issues \#61 a \#79).}\\
\toprule
\textbf{Fecha} & \textbf{Marcelo (motor e IA)} & \textbf{Abner (plataforma y seguridad)}\\
\midrule
\endfirsthead
\toprule
\textbf{Fecha} & \textbf{Marcelo (motor e IA)} & \textbf{Abner (plataforma y seguridad)}\\
\midrule
\endhead
16/11 lun & Prompts de generación con contexto GraphRAG. & Plantilla Copier: módulo, entidades Prisma y DTO.\\
17/11 mar & Puntos de extensión y lógica de negocio con el LLM. & Plantilla: servicio, controlador y decoradores Swagger.\\
18/11 mié & Lista de dependencias permitidas y verificación. & Plantilla: guard Keycloak, ValidationPipe, Helmet y filtro de errores.\\
19/11 jue & Generación de pruebas Jest con el LLM. & Plantilla: Dockerfile, Jest base y configuración por entorno.\\
23/11 lun & Generador: recorrer \texttt{arquitectura.json} e invocar Copier y el LLM. & Parámetros de la plantilla por servicio y validación de la estructura.\\
24/11 mar & Verificación con \texttt{tsc -{}-noEmit} y reintento con el LLM. & Gate 1 sobre los servicios generados.\\
25/11 mié & Trazabilidad requisito → componente → archivo en Neo4j. & Ajustes de la plantilla según los hallazgos del Gate 1.\\
26/11 jue & Fases 1--4 con los documentos 1 y 2; compilación al primer intento. & Documento 3 y revisión de los hallazgos del Gate 1.\\
\rowcolor{azulclaro} 27/11 vie & Informe del sprint 4 y métricas iniciales. & Demostración y planificación del sprint 5.\\
\bottomrule
\end{longtable}}

\appendix
% ======================================================================
\chapter{Evidencias}
% ======================================================================

\begin{table}[H]
\centering
\caption{Dónde comprobar cada resultado del sprint.}
\begin{tabularx}{\textwidth}{@{}P{4.0cm}>{\raggedright\arraybackslash\small}X@{}}
\toprule
\textbf{Evidencia} & \textbf{Ubicación}\\
\midrule
Ejecución de CI al cierre & \url{https://github.com/PracticasESPE/Framework_Desarrollo_Seguro/actions/runs/43660815603}\\
Dataset de evaluación & \texttt{evaluation/datasets/} (tres documentos, cuatro formatos y su referencia)\\
Fase 1 & \texttt{orchestrator/ingest/}; evaluación en \texttt{evaluation/fase1/}\\
Fase 2 y ciclo LLM ↔ OPA & \texttt{orchestrator/architecture/}, \texttt{policies/arquitectura/}; evaluación en \texttt{evaluation/fase2/}\\
ADR del documento 1 & \texttt{docs/adr/} (con sus entradas en \texttt{docs/adr/entrada/})\\
Diagrama C4 del documento 1 & \texttt{docs/arquitectura.md}\\
Dashboard y API & \texttt{dashboard/} y \texttt{orchestrator/api.py}; documentación en \texttt{/docs} de la API\\
Gestión de los LLM & \texttt{orchestrator/llm/README.md} (cuota, caché, respaldo y comparación)\\
Esquemas JSON & \texttt{docs/esquemas/modelo\_requisitos.schema.json} y \texttt{arquitectura.schema.json}\\
Guía de la demostración & \texttt{docs/sprints/sprint3-demostracion.md}\\
Planificación del sprint 3 & \texttt{docs/sprints/sprint3-planificacion.md}\\
\bottomrule
\end{tabularx}
\end{table}

\section*{Cómo reproducir las fases 1 y 2 en local}
Con Neo4j en marcha, OPA instalado y la clave del LLM en \texttt{.env}, desde la raíz del
repositorio (en Windows con tildes en la ruta, \texttt{python -m orchestrator.cli} en lugar de
\texttt{framework}):

\begin{tcolorbox}[colback=white,colframe=grisfw,boxrule=0.4pt,arc=1mm,left=2mm]
\small\ttfamily
uv run framework requisitos clasificar evaluation/datasets/01\_tienda\_en\_linea/requisitos.docx\\
uv run python -m evaluation.fase1.evaluar salida/requisitos/requisitos-docx.modelo\_requisitos.json evaluation/datasets/01\_tienda\_en\_linea\\
uv run framework arquitectura proponer salida/requisitos/requisitos-docx.modelo\_requisitos.json\\
uv run framework arquitectura adr salida/arquitectura/requisitos-docx.arquitectura.json -{}-requisitos salida/requisitos/requisitos-docx.modelo\_requisitos.json -{}-destino salida/adr\\
uv run framework arquitectura diagrama salida/arquitectura/requisitos-docx.arquitectura.json -{}-requisitos salida/requisitos/requisitos-docx.modelo\_requisitos.json -{}-salida salida/arquitectura.md -{}-adr salida/adr
\end{tcolorbox}

\end{document}
